\section{因子分析与概率 PCA：把低秩变成生成模型}
\label{sec:13-Machine-Learning/04-Dimensionality-Reduction/06}

一份 40 题的员工问卷，人事想要三个维度的分数。你跑了 PCA，前三个成分解释了 \(62\%\) 的方差，交上去。回来三个问题：第二个成分上那些正负号代表什么？为什么下一季度加了两道题，三个成分的载荷全变了？还有——问卷里有几道题公认答得很随意，测量噪声远大于其他题，PCA 有没有把这件事考虑进去？三个问题 PCA 都答不上来，因为它压根没有``噪声''这个概念。它只有方差。

\subsection{相关变量背后可能只有少数共同原因}

先把直觉说清楚。问卷里的题目不是彼此独立的：一个人对工作满意，多数正面题目会一起给高分；心生去意时，多数题目一起往下走。看数据的相关矩阵，一大片正相关。你猜背后只有少数几个共同的东西在驱动，比如对直属主管的评价、对薪酬的评价、对成长空间的评价；剩下的波动是每道题自己的措辞效应和答题随机性。

PCA 的做法是找一个低维子空间使重构误差最小。它抓住了主要变化方向，但它对所有坐标一视同仁——隐含假设是各方向的噪声同等重要。于是那几道答得随意的题，因为方差大，会把主方向往自己那边拽。想把``大家一起变''和``各变各的''分开，需要在模型里给后者留一个位置。

因子分析换了一条路：不直接找子空间，而是写下一个生成过程。潜变量 \(z\) 是那几个共同原因，观测 \(x\) 是它们的线性映像加上各自独立的噪声，

\[
z\sim\mathcal N(0,I_k),
\qquad
x=\mu+\Lambda z+\varepsilon,
\qquad
\varepsilon\sim\mathcal N(0,\Psi),
\]

其中 \(\Lambda\in\R^{p\times k}\) 是因子载荷矩阵，每一列描述一个潜因子对所有题目的影响分布；\(\Psi\) 是正对角矩阵，第 \(i\) 个对角元是第 \(i\) 道题的特有方差，也就是共同因子解释不了的那部分。

\subsection{把低秩写进生成过程，协方差就自动裂成两块}

把 \(z\) 积分掉，观测的边缘分布是

\[
x\sim\mathcal N\bigl(\mu,\ \Lambda\Lambda^{\trans}+\Psi\bigr).
\]

这一行是整个模型的全部内容。协方差被拆成两项：\(\Lambda\Lambda^{\trans}\) 秩不超过 \(k\)，承担共同变化；\(\Psi\) 是对角的，承担各自独立的波动。

\begin{center}
\begin{tikzpicture}[>=stealth,scale=1.0]
% Sigma
\draw[black!70,fill=black!10] (0,0) rectangle (2.0,2.0);
\node[font=\small] at (1.0,1.0) {\(\Sigma\)};
\node[font=\scriptsize,anchor=north] at (1.0,-0.15) {\(p\times p\)};
\node[font=\large] at (2.45,1.0) {\(=\)};
% 低秩部分
\draw[black!70,fill=black!22] (2.9,0) rectangle (3.45,2.0);
\node[font=\scriptsize] at (3.17,1.0) {\(\Lambda\)};
\draw[black!70,fill=black!22] (3.6,1.45) rectangle (5.6,2.0);
\node[font=\scriptsize] at (4.6,1.72) {\(\Lambda^{\trans}\)};
\node[font=\scriptsize,anchor=north] at (4.25,-0.15) {秩 \(\le k\)};
\node[font=\large] at (5.95,1.0) {\(+\)};
% Psi
\draw[black!70] (6.4,0) rectangle (8.4,2.0);
\foreach \i in {0,1,2,3,4,5,6,7}
  \fill[black!55] ({6.4+\i*0.25},{2.0-\i*0.25-0.25}) rectangle ({6.4+\i*0.25+0.25},{2.0-\i*0.25});
\node[font=\scriptsize,anchor=north] at (7.4,-0.15) {\(\Psi\) 只有对角};
\end{tikzpicture}
\end{center}

\emph{因子模型对协方差的全部主张：非对角的相关必须由那 \(k\) 列载荷解释，剩下的只许留在对角线上。}

图里最右边那块方阵只有对角线有值，这不是画图偷懒，而是模型最强的假设：\textbf{给定潜因子之后，各题目条件独立}。所有的题间相关都必须由 \(\Lambda\Lambda^{\trans}\) 生产出来。这条假设也给了诊断的入口——拟合之后算残差协方差 \(S-\widehat\Lambda\widehat\Lambda^{\trans}-\widehat\Psi\)，若非对角上还剩下成片的结构，说明因子数不够，或者有一对题目共享了模型之外的东西（比如措辞几乎一样）。

参数用最大似然估计，见\hyperref[sec:11-Mathematical-Statistics/02-Point-Estimation/03]{``最大似然选择最能解释数据的参数''}一节；实现上通常用 EM，把 \(z\) 当作缺失数据，E 步算 \(p(z\given x)\) 的一阶二阶矩，M 步更新 \(\Lambda\) 与 \(\Psi\)。

\subsection{PCA 是噪声消失时的极限，但两者回答的不是一个问题}

概率 PCA 是因子分析的一个特例：把 \(\Psi\) 收紧成 \(\sigma^2 I\)，即所有变量的噪声方差相同。这时最大似然解有闭式——\(\Lambda\) 的列空间恰好由样本协方差最大的 \(k\) 个特征方向张成，\(\sigma^2\) 的估计是剩下 \(p-k\) 个特征值的平均。让 \(\sigma^2\to 0\)，潜变量的后验均值退化成普通 PCA 的正交投影。

所以 PCA 是 PPCA 在噪声消失时的极限。但``是极限''不等于``是同一件事''。PCA 给的是一个确定性的坐标变换；PPCA 给的是一个完整的联合分布 \(p(x,z)\)，于是多出来一整套东西：潜变量的后验 \(p(z\given x)\) 带方差，样本少或噪声大时后验会自己变宽；边缘似然 \(p(x)\) 可以拿来比较不同的 \(k\)，也可以拿来做异常检测；缺失值不需要预先填补，直接从 \(p(x_{\text{miss}}\given x_{\text{obs}})\) 推断即可。这些能力不是给 PCA 的公式换了个概率名字，而是由``写下了一个分布''本身带来的。换来这些的前提也明确：你得相信那个各向同性高斯噪声，以及潜变量的高斯先验。

回到开头的第三个问题。若某几道题的测量噪声明显更大，PPCA 的 \(\Psi=\sigma^2I\) 无法表达这一点，多出来的噪声会被强行匀进共同因子，污染载荷；经典因子分析可以给这几道题更大的 \(\psi_i\)，让它们对公共因子少说话。这是选 FA 还是选 PPCA 的实质分歧，不是术语偏好。

\subsection{载荷的列没有唯一方向，可辨认的只有子空间}

解出 \(\Lambda\) 之后，很自然会逐列去读：这一列像是对主管的评价，那一列像是对薪酬的评价。先慢一步。对任意正交矩阵 \(R\)，

\[
\Lambda\Lambda^{\trans} = (\Lambda R)(\Lambda R)^{\trans},
\]

似然只依赖 \(\Lambda\Lambda^{\trans}\)。旋转载荷，协方差不变、似然不变、拟合优度一个小数点都不动，而每一列的含义可以面目全非。数据能确定的是这 \(k\) 列张成的子空间，不是列本身。这和上一节矩阵补全里 \((U,V)\) 可以同时旋转、以及第二节里奇异值接近时奇异向量不可解释，是同一类现象的三次出现。

于是所谓的因子旋转——varimax 让每列尽量集中在少数题目上，oblimin 允许因子彼此相关——是人在子空间里挑一个便于讲述的基。挑得好，载荷表更容易和领域知识对齐；但这份``容易''是你加进去的偏好，不是数据给出的证据。开头那个``第二个成分的正负号是什么意思''，正确的回答是：单看一列，没有意思。

尺度和符号也要靠约定钉死。若允许 \(\Cov(z)=\Sigma_z\) 自由，则 \((\Lambda,\Sigma_z)\) 与 \((\Lambda T,\ T^{-1}\Sigma_z T^{-\trans})\) 给出同一个 \(\Lambda\Sigma_z\Lambda^{\trans}\)，载荷的尺度和潜变量的尺度可以互相补偿。设 \(\Cov(z)=I\) 就是在拿掉这层冗余。可辨认性约束不是形式上的细枝末节，它划定了哪些参数陈述有统计含义、哪些只是约定的影子。

\subsection{特有方差给了灵活性，也给了退化的机会}

让每个变量有自己的 \(\psi_i\)，参数因此变多，似然曲面也变复杂。小样本下 \(\Lambda\) 与 \(\Psi\) 未必分得开：数据不够时，说不清某个方向上的方差是共同因子造的还是特有噪声造的。更具体的病症是 Heywood 情形——某个 \(\psi_i\) 被估到零甚至负值，字面意思是这道题被公共因子完美解释、没有任何自身波动。真实数据里这几乎不可能，它通常是因子数取多了、样本太小或模型误设的信号，而不是一个好消息。

因子数的选择同样不能只看碎石图。第一节说过，高维下样本协方差的谱本身就会散开，肉眼找到的肘部可能纯属采样涨落，判断需要 Marchenko--Pastur 那样的零模型作基准，见\hyperref[ch:094]{``随机矩阵理论：高维协方差的谱失真''}一章。在因子模型里还有更直接的手段：比较不同 \(k\) 的留出数据似然，看残差协方差是否已无结构，用 bootstrap 检查载荷子空间在重抽样下是否稳定。一张漂亮的二维载荷图撑不起结论，这几项检查才撑得起。

\subsection{生成模型不会顺带赋予潜变量现实身份}

最后一条边界，也是最容易越过的一条。因子分析做的事情是：给定一个协方差结构，找一个最简洁的生成描述。它没有能力判断潜因子是否真的存在、是否可被干预、是否在不同人群里指的是同一个东西。

数据经过筛选（只调查了某类员工）、量表设计（措辞诱导了特定回答）、或预处理（标准化、去均值）之后，因子结构会随之改变。不同群体估出的载荷也可能不同，这就是测量不变性问题：同一份问卷在两个群体里，测的可能不是同一个构念，此时把因子得分跨群比较是没有依据的。把因子分析用在问卷分数或风险评分上时，跨群的载荷一致性应当先验证，再使用。

低秩告诉你许多变量可以在少数维度上共同变化，概率模型告诉你共同变化与噪声如何在似然框架里分开。两者都没有回答那些方向对应现实中的什么。若想让方向本身而不只是子空间变得可辨认，必须往模型里再加一条结构假设——下一节的两种加法，一种赌非高斯，一种赌两组变量之间的相关。
